\chapter{Myomectomy}

\section*{What Is This Surgery?}
A myomectomy is a surgical procedure aimed at the excision of uterine fibroids, also known as leiomyomas, while preserving the uterus. This intervention is primarily indicated for women who experience significant symptoms attributable to fibroids, such as heavy menstrual bleeding, pelvic pain, and pressure, and who wish to maintain their fertility potential or avoid a hysterectomy. The procedure involves the careful removal of benign tumors from the myometrial wall, followed by meticulous reconstruction of the uterine musculature to restore its integrity and function. Myomectomy can be performed through various approaches, including open abdominal, laparoscopic, hysteroscopic, or robotic techniques, each selected based on the size, number, and anatomical location of the fibroids, as well as patient-specific factors and surgeon expertise. The clinical significance of myomectomy lies in its ability to alleviate debilitating symptoms, improve quality of life, and enhance reproductive prospects for women suffering from fibroid-related complications.

\section*{Alternative Names}
\begin{itemize}
  \item Uterine Fibroid Removal Surgery
  \item Leiomyoma Excision
  \item Laparoscopic Myomectomy
  \item Hysteroscopic Myomectomy
  \item Abdominal Myomectomy
  \item Robotic Myomectomy
\end{itemize}

\section*{Relevant Surgical Anatomy}
A comprehensive understanding of pelvic anatomy is crucial for the safe and effective performance of myomectomy. The uterus is a pear-shaped muscular organ located in the true pelvis, positioned between the bladder anteriorly and the rectum posteriorly. It consists of three main parts: the fundus, body (corpus), and cervix. The uterine wall is composed of three layers: the outer serosa (perimetrium), the thick muscular myometrium where fibroids typically arise, and the inner endometrium. 

The arterial supply to the uterus is primarily derived from the uterine arteries, which are branches of the internal iliac arteries. These arteries run within the broad ligament and cross anterior to the ureters at the level of the cervix. The ovarian arteries, which are direct branches of the aorta, supply the ovaries and contribute collateral circulation to the uterine fundus through anastomoses with the uterine arteries. Venous drainage parallels the arterial supply, forming extensive plexuses that drain into the internal iliac veins and ovarian veins. 

Nerve supply is provided by the inferior hypogastric plexus, which offers sympathetic and parasympathetic innervation. Lymphatic drainage follows the major blood vessels, primarily to the internal and external iliac nodes and obturator nodes. Surgical landmarks include the ureters, which must be identified and protected during dissection, particularly when fibroids are large and distort normal anatomy. The "pseudocapsule" of a fibroid, a layer of compressed myometrial tissue and connective tissue, serves as a critical surgical plane for enucleation.

\section*{Surgical Principle \& Rationale}
The fundamental principle of myomectomy is the enucleation of uterine fibroids from the myometrial wall while preserving the surrounding healthy uterine tissue. This involves identifying the fibroid, incising the overlying serosa and myometrium, and carefully dissecting the fibroid from its pseudocapsule. The pseudocapsule often provides a natural plane for dissection, facilitating removal and minimizing damage to the adjacent myometrium. 

A critical aspect of the procedure is achieving meticulous hemostasis during and after fibroid removal, often through the use of vasoconstrictive agents (e.g., dilute vasopressin), electrosurgery, or sutures, to prevent excessive blood loss. Following fibroid removal, the defect in the uterine wall must be meticulously repaired in multiple layers to restore the anatomical integrity and functional strength of the uterus. This layered closure is crucial for minimizing postoperative adhesion formation, preventing uterine dehiscence, and ensuring the uterus can safely withstand a future pregnancy. The specific surgical approach dictates the method of access and removal: hysteroscopic myomectomy is used for submucosal fibroids protruding into the uterine cavity, laparoscopic and robotic myomectomies are employed for subserosal and intramural fibroids, and open abdominal myomectomy is reserved for very large, numerous, or deeply embedded fibroids. Regardless of the approach, the overarching goal is to eliminate the symptomatic fibroids while safeguarding the uterus for its reproductive and structural roles.

\section*{History \& Surgical Evolution}
The concept of removing uterine fibroids while preserving the uterus has evolved significantly over centuries. Early attempts at myomectomy were fraught with high morbidity and mortality due to challenges in controlling hemorrhage and managing infection. The first documented successful abdominal myomectomy is often attributed to Dr. Alexander J.C. Skene in 1898, who performed the procedure on a patient with a large uterine fibroid, demonstrating the feasibility of uterine preservation. Prior to this, hysterectomy was the standard treatment for symptomatic fibroids.

The early 20th century saw gradual improvements in surgical techniques, anesthesia, and aseptic practices, making abdominal myomectomy a safer and more viable option. However, it remained a major open surgery with significant recovery times. The latter half of the 20th century ushered in a revolution in gynecological surgery with the advent of minimally invasive techniques. Hysteroscopic myomectomy, allowing removal of submucosal fibroids through the cervix, was pioneered in the 1970s and 1980s. Shortly thereafter, laparoscopic myomectomy emerged, offering a less invasive approach for intramural and subserosal fibroids, with early successes reported by surgeons like Harry Reich in the late 1980s and early 1990s. The 21st century has further advanced the field with the introduction of robotic-assisted laparoscopic myomectomy, which provides enhanced dexterity, 3D visualization, and improved ergonomics for complex cases, further refining the precision and safety of uterine fibroid removal.

\section*{Clinical Indications}
Myomectomy is indicated for women experiencing significant symptoms due to uterine fibroids, particularly when fertility preservation is a priority. Specific indications include:

\begin{itemize}
  \item Heavy or prolonged menstrual bleeding (menorrhagia) leading to anemia and fatigue, often refractory to medical management.
  \item Severe pelvic pain, pressure, or discomfort, which may be chronic or acute, impacting daily activities.
  \item Bulk-related symptoms such as urinary frequency, urgency, constipation, or a feeling of pelvic fullness due to fibroid compression on adjacent organs.
  \item Infertility or recurrent pregnancy loss where fibroids are identified as a likely contributing factor, especially submucosal or large intramural fibroids distorting the uterine cavity.
  \item Rapid growth of a uterine fibroid, which may raise concern for malignancy, although leiomyosarcoma is rare and typically not diagnosed preoperatively.
  \item Fibroids causing dyspareunia (painful intercourse) or other sexual dysfunction due to their size or location.
  \item Torsion of a pedunculated subserosal fibroid, causing acute abdominal pain requiring urgent surgical intervention.
  \item Anemia refractory to medical management, directly attributable to fibroid-induced bleeding, necessitating surgical intervention to control blood loss.
\end{itemize}

\section*{Clinical Positioning}
Myomectomy is generally considered a primary surgical treatment for symptomatic uterine fibroids, especially in women who desire to preserve their uterus and potential for future fertility. It is often the first-line surgical option when medical management (such as hormonal therapies or NSAIDs) has failed to adequately control symptoms, or when fibroids are too large or numerous for non-surgical interventions like uterine artery embolization (UAE) to be effective or appropriate. 

In cases where fertility is a concern, myomectomy is preferred over hysterectomy, which is a definitive treatment but eliminates reproductive capacity. While myomectomy is primary for fertility-sparing situations, it can be considered a secondary or salvage procedure in certain contexts. For instance, if a woman initially undergoes UAE but experiences persistent symptoms or fibroid regrowth, myomectomy might be considered as a secondary intervention. Similarly, in cases of failed medical management or when a patient develops new or worsening symptoms after a period of observation, myomectomy becomes the next step. It is also an alternative to hysterectomy for women who do not desire fertility but wish to retain their uterus for personal or psychological reasons, making it a primary choice in that specific scenario as well. The choice of approach (hysteroscopic, laparoscopic, robotic, or open) is determined by fibroid characteristics and patient factors, positioning myomectomy as a versatile primary intervention.

\section*{Surgical Technique \& Procedure}
Myomectomy is performed under general anesthesia, ensuring the patient is completely unconscious and pain-free throughout the procedure. Patient positioning varies by approach: supine for abdominal or robotic myomectomy, often with Trendelenburg for laparoscopy, and lithotomy for hysteroscopy.

The key operative steps, with variations based on approach, typically include:

\begin{enumerate}
  \item \textbf{Access and Visualization:} For abdominal approaches (open, laparoscopic, robotic), an incision is made (Pfannenstiel or midline for open; umbilical and accessory ports for laparoscopic/robotic) and the abdomen is insufflated with carbon dioxide to create working space. For hysteroscopic myomectomy, the cervix is dilated, and a resectoscope is inserted into the uterine cavity, with continuous infusion of distending media.
  
  \item \textbf{Uterine Manipulation and Hemostasis:} The uterus is manipulated to expose the fibroids. For abdominal approaches, a dilute vasopressin solution may be injected into the myometrium surrounding the fibroid to induce vasoconstriction and minimize blood loss. Tourniquets or temporary clamping of uterine arteries may also be employed in open cases.
  
  \item \textbf{Fibroid Incision and Enucleation:} An incision is made over the fibroid, either directly through the uterine serosa and myometrium (abdominal approaches) or by shaving/excising tissue within the uterine cavity (hysteroscopic). The fibroid is then carefully dissected and enucleated from its pseudocapsule using specialized instruments, preserving surrounding healthy myometrium.
  
  \item \textbf{Uterine Wall Repair:} After fibroid removal, the defect in the uterine wall is meticulously repaired. For abdominal approaches, this involves a multi-layered closure of the myometrium and serosa using absorbable sutures to restore anatomical integrity and prevent uterine dehiscence. For hysteroscopic cases, no repair is typically needed as tissue is shaved from the cavity.
  
  \item \textbf{Fibroid Extraction:} Removed fibroids from abdominal approaches are extracted. Smaller fibroids may be removed through an enlarged port incision or morcellated (cut into smaller pieces) within a containment bag to prevent tissue dissemination. For hysteroscopic cases, tissue fragments are extracted from the uterus via the resectoscope.
  
  \item \textbf{Hemostasis and Adhesion Prevention:} The surgical site is thoroughly inspected for hemostasis. Adhesion barrier agents (e.g., oxidized regenerated cellulose, hyaluronic acid-carboxymethylcellulose) may be applied to the uterine incision site to minimize postoperative adhesion formation, particularly in cases of extensive dissection.
  
  \item \textbf{Closure:} For abdominal approaches, the abdominal incision(s) are closed in layers. Drains are rarely used but may be considered in cases of extensive dissection or significant oozing. For hysteroscopic cases, the resectoscope is removed, and the cervix naturally closes.
\end{enumerate}

\section*{Preoperative Workup \& Preparation}
Thorough preoperative preparation is essential for a safe and successful myomectomy. This typically includes a comprehensive medical and cardiac evaluation to assess overall health and fitness for general anesthesia, often involving an anesthesiologist and, if necessary, a cardiologist.

\begin{itemize}
  \item \textbf{Diagnostic Imaging:} Pelvic ultrasound is standard, but magnetic resonance imaging (MRI) is often preferred to precisely map the size, number, and location of fibroids, assess their relationship to the uterine cavity and serosa, and rule out other pathologies. Hysterosalpingography or sonohysterography may be used to evaluate the uterine cavity for submucosal fibroids or other abnormalities.
  
  \item \textbf{Laboratory Tests:} A complete blood count (CBC) is performed to check for anemia, blood type and cross-match are obtained in case of transfusion, coagulation studies are done to assess bleeding risk, and a basic metabolic panel evaluates kidney function and electrolytes.
  
  \item \textbf{Medication Adjustments:} Patients may need to stop certain medications, such as blood thinners (e.g., aspirin, NSAIDs, warfarin, novel oral anticoagulants), several days to a week before surgery to minimize bleeding risk.
  
  \item \textbf{Preoperative Medications:}
    \begin{itemize}
      \item \textit{GnRH Agonists:} Often prescribed for 3-6 months prior to surgery to shrink fibroids, reduce uterine vascularity, and improve anemia, especially for large fibroids.
      \item \textit{Iron Supplementation:} For anemic patients, to optimize hemoglobin levels before surgery.
      \item \textit{Prophylactic Antibiotics:} Administered intravenously shortly before incision to reduce the risk of surgical site infection.
    \end{itemize}
  
  \item \textbf{NPO Status:} Patients are instructed to be NPO (nil per os) for a specified period (typically 6-8 hours) before surgery to prevent aspiration during anesthesia.
  
  \item \textbf{Bowel Preparation:} May be required for laparoscopic or robotic cases if there is a high risk of bowel injury or if extensive dissection is anticipated, or if the fibroids are very large.
  
  \item \textbf{Informed Consent:} A detailed discussion with the surgeon regarding the procedure, potential risks, benefits, alternatives, and the possibility of conversion to hysterectomy is crucial.
\end{itemize}

\section*{Contraindications \& Precautions}
\textbf{Absolute Contraindications:}

\begin{itemize}
  \item \textbf{Suspected Uterine Malignancy (Leiomyosarcoma):} If preoperative imaging or clinical findings strongly suggest leiomyosarcoma, myomectomy is contraindicated due to the risk of tumor dissemination. Hysterectomy with appropriate oncologic staging is indicated.
  
  \item \textbf{Active Pelvic Infection:} Acute pelvic inflammatory disease or other active pelvic infections must be treated and resolved before elective myomectomy to prevent sepsis and other infectious complications.
  
  \item \textbf{Severe Uncontrolled Medical Comorbidities:} Unstable cardiac disease, severe pulmonary disease, uncontrolled diabetes, or other life-threatening conditions that make general anesthesia and major surgery excessively risky.
  
  \item \textbf{Pregnancy:} Myomectomy is generally contraindicated during pregnancy, except in rare, emergent circumstances such as torsion of a pedunculated fibroid, due to high risks of miscarriage or preterm labor.
\end{itemize}

\textbf{Relative Contraindications \& Precautions:}

\begin{itemize}
  \item \textbf{Extremely Large or Numerous Fibroids:} While not an absolute contraindication, very large uteri (e.g., >20 weeks size) or an excessive number of fibroids (e.g., >10-15) can make myomectomy technically challenging, increase blood loss, and prolong operative time, potentially increasing the risk of complications or requiring conversion to hysterectomy.
  
  \item \textbf{Desire for Immediate Conception:} After extensive myomectomy, particularly open abdominal, a healing period of 3-6 months is typically recommended before attempting pregnancy to allow the uterine scar to strengthen, reducing the risk of uterine rupture.
  
  \item \textbf{Previous Extensive Uterine Surgery:} Prior uterine surgeries, especially those involving significant myometrial incisions (e.g., classical cesarean section), can increase the risk of adhesions and make subsequent myomectomy more difficult and risky.
  
  \item \textbf{Patient Age and Parity:} For women nearing menopause or those who have completed childbearing and do not wish to preserve their uterus, hysterectomy may be a more definitive and less complex option, with a lower risk of fibroid recurrence.
  
  \item \textbf{Risk of Fibroid Recurrence:} Patients should be counseled that myomectomy removes existing fibroids but does not prevent the development of new fibroids, and recurrence rates vary depending on age and number of fibroids removed.
  
  \item \textbf{Need for Cesarean Delivery:} After a transmural uterine incision (e.g., for intramural fibroids), future pregnancies will typically require delivery by cesarean section to prevent uterine rupture during labor.
  
  \item \textbf{Potential for Conversion to Hysterectomy:} Patients must understand that in rare cases of uncontrollable bleeding or unexpected findings (e.g., suspected malignancy), conversion to hysterectomy may be necessary to ensure patient safety.
\end{itemize}

\section*{Complications \& Risks}
As with any surgical procedure, myomectomy carries potential risks, which can be broadly categorized:

\begin{itemize}
  \item \textbf{General Surgical Complications:}
    \begin{itemize}
      \item \textit{Bleeding:} Intraoperative hemorrhage, potentially requiring blood transfusion or, in severe cases, conversion to hysterectomy. Postoperative hematoma formation is also a risk.
      \item \textit{Infection:} Surgical site infection (wound infection), pelvic infection (e.g., endometritis, abscess), or urinary tract infection.
      \item \textit{Anesthesia Complications:} Reactions to anesthetic agents, respiratory problems (e.g., pneumonia, atelectasis), cardiac events (e.g., arrhythmia, myocardial infarction), or nerve injury from prolonged positioning.
      \item \textit{Thromboembolic Events:} Deep vein thrombosis (DVT) or pulmonary embolism (PE), though rare, can be life-threatening. Prophylactic measures are routinely employed.
    \end{itemize}
  
  \item \textbf{Procedure-Specific Risks:}
    \begin{itemize}
      \item \textit{Uterine Perforation:} Especially a risk during hysteroscopic myomectomy, potentially requiring laparoscopic repair or conversion to open surgery.
      \item \textit{Adhesion Formation:} Scar tissue formation within the pelvis or around the uterus, which can lead to chronic pelvic pain, bowel obstruction, or infertility.
      \item \textit{Injury to Adjacent Organs:} Accidental damage to the bladder, bowel, ureters, or major blood vessels, potentially requiring further surgical repair.
      \item \textit{Fibroid Recurrence:} Myomectomy removes existing fibroids but does not prevent new ones from growing. Recurrence rates vary but can be significant over time, potentially necessitating further intervention.
      \item \textit{Incomplete Fibroid Removal:} Small, undetected fibroids may be left behind, or a large fibroid may not be entirely removed, leading to persistent symptoms.
      \item \textit{Uterine Rupture in Future Pregnancy:} A rare but serious complication, particularly after extensive transmural myomectomy, necessitating a planned cesarean delivery for subsequent pregnancies.
      \item \textit{Conversion to Hysterectomy:} In cases of uncontrollable bleeding, unexpected malignancy, or unforeseen complications, the surgeon may need to perform a hysterectomy to ensure patient safety.
      \item \textit{Fluid Overload (Hysteroscopic Myomectomy):} Excessive absorption of distending media can lead to hyponatremia, pulmonary edema, and cardiac compromise.
      \item \textit{Nerve Injury:} Due to prolonged positioning or pressure during surgery, though uncommon, can result in temporary or permanent sensory or motor deficits.
    \end{itemize}
\end{itemize}

\section*{Postoperative Recovery Timeline}
Immediately after myomectomy, the patient is transferred to the Post-Anesthesia Care Unit (PACU) for close monitoring as they recover from anesthesia. Pain management is initiated, typically with intravenous analgesics, and vital signs are continuously observed. The length of hospital stay varies significantly by surgical approach: hysteroscopic myomectomy is often an outpatient procedure or requires an overnight stay, laparoscopic/robotic myomectomy usually involves 1-2 nights, and open abdominal myomectomy typically requires 2-4 nights.

During the first 24-48 hours, patients are encouraged to ambulate early to prevent complications like DVT and promote bowel function. Pain is managed with oral or intravenous medications, and anti-nausea medication may be given. Diet is gradually advanced from clear liquids to solids as tolerated. For abdominal incisions, wound care instructions are provided. Over the first 1-2 weeks, patients will experience varying degrees of pain, fatigue, and vaginal spotting or discharge. Activity restrictions are crucial: heavy lifting, strenuous exercise, and sexual intercourse are typically avoided for 4-6 weeks to allow the uterus to heal. Full recovery and return to normal activities, including work, can take 2-4 weeks for minimally invasive approaches and 4-8 weeks for open abdominal surgery. Long-term, patients are advised on the appropriate waiting period before attempting pregnancy, usually 3-6 months, to ensure adequate uterine scar healing and minimize the risk of uterine rupture.

\section*{Surgical Findings \& Pathology}
During myomectomy, the surgeon documents several key findings, including the number, size, and precise location of all identified fibroids (e.g., subserosal, intramural, submucosal, pedunculated). The presence of any associated pelvic adhesions, endometriosis, or other pelvic pathology is also noted. The appearance of the fibroids, such as areas of degeneration (hyaline, cystic, myxoid, red degeneration), is observed. All removed fibroid tissue is sent to a pathology laboratory for microscopic examination. The pathology report is crucial for confirming the diagnosis of benign leiomyomas, ruling out the rare but serious possibility of leiomyosarcoma. The report will detail the histological characteristics of the fibroids, including cellularity, mitotic activity, and any degenerative changes. This pathological confirmation is vital for patient reassurance and guiding future management.

\section*{Expected Postoperative Course}
A normal and successful postoperative course following myomectomy is characterized by a progressive resolution of symptoms and a return to baseline health. Patients can expect initial incisional pain, which should gradually decrease over days to weeks, managed effectively with prescribed analgesics. Vaginal spotting or light bleeding is common for several days to a few weeks. Bowel function typically returns within 24-48 hours, and patients should tolerate a regular diet. Wound healing progresses without signs of infection, and any sutures or staples are removed at the first follow-up visit. Fatigue is common initially but improves with rest and gradual activity. The primary goal is the significant reduction or complete alleviation of preoperative symptoms such as heavy menstrual bleeding, pelvic pain, and pressure, leading to an improved quality of life. For those with fertility goals, the uterus is expected to heal robustly, preparing for future conception.

\section*{Postoperative Care \& Follow-up}
Postoperative care and follow-up are crucial for monitoring recovery, addressing any concerns, and ensuring optimal long-term outcomes.

\begin{itemize}
  \item \textbf{Postoperative Visits:} Typically, a first follow-up visit is scheduled 1-2 weeks after surgery to check wound healing, discuss the pathology report, and remove any non-absorbable sutures or staples. A subsequent visit may occur at 6-8 weeks to assess overall recovery and discuss long-term management.
  
  \item \textbf{Imaging:} A follow-up ultrasound may be performed several months after surgery to assess the healed uterine cavity, confirm the absence of residual fibroids, and evaluate for any new fibroid growth, especially if symptoms recur.
  
  \item \textbf{Activity Resumption:} Gradual resumption of normal activities, including exercise, sexual activity, and return to work, is advised based on the surgical approach and individual recovery.
  
  \item \textbf{Fertility Counseling:} For women desiring future pregnancy, counseling on the optimal timing for conception (typically 3-6 months post-surgery) and the necessity of a planned cesarean delivery if a transmural uterine incision was made is provided.
  
  \item \textbf{Warning Signs:} Patients are educated on warning signs to report immediately, such as fever, increasing abdominal pain, heavy vaginal bleeding (more than a menstrual period), foul-smelling vaginal discharge, or signs of wound infection (redness, swelling, pus).
  
  \item \textbf{Long-term Management:} Discussion about the risk of fibroid recurrence and potential future management options, including lifestyle modifications, medical therapies, or further surgical intervention, is part of ongoing care.
\end{itemize}

Adherence to these follow-up recommendations is vital for sustained health and well-being.

\section*{Epidemiology}
Uterine fibroids are the most common benign tumors in women, affecting an estimated 70-80\% of women by age 50. However, only a subset of these women develop symptoms severe enough to warrant intervention. Myomectomy is a frequently performed gynecological procedure, particularly among women in their reproductive years. The incidence of myomectomy is highest in women aged 30-45, with a peak in the late 30s and early 40s, reflecting the typical age range for symptomatic fibroids and fertility preservation desires. African American women have a higher incidence of fibroids, develop them at an earlier age, and tend to have larger and more numerous fibroids, leading to a disproportionately higher rate of myomectomy compared to Caucasian women. The most common indications for myomectomy are heavy menstrual bleeding and pelvic pain, followed by infertility. While myomectomy is generally safe, the overall morbidity rate is low, with serious complications occurring in less than 5\% of cases. Mortality is exceedingly rare, estimated at less than 0.1\%.

\section*{Outcomes \& Prognosis}
The outcomes of myomectomy are generally favorable, particularly for symptom relief and fertility preservation. Typical success rates for alleviating symptoms like heavy menstrual bleeding and pelvic pain range from 80-90\%. For women with infertility attributed to fibroids, myomectomy can significantly improve pregnancy rates, with studies reporting conception rates between 40-60\% after surgery, depending on the extent of fibroid disease and other fertility factors. Functional outcomes, such as improved quality of life and reduced need for pain medication, are also high. The American College of Obstetricians and Gynecologists (ACOG) guidelines support myomectomy as an effective treatment for symptomatic fibroids in women desiring uterine preservation.

However, fibroid recurrence is a significant prognostic factor. While myomectomy removes existing fibroids, it does not prevent the growth of new ones. Recurrence rates vary widely, from 10-50\% over 5-10 years, influenced by factors such as the patient's age (younger women have higher recurrence), the number and size of fibroids removed, and genetic predisposition. Approximately 10-25\% of women who undergo myomectomy may eventually require a repeat myomectomy or a hysterectomy due to recurrent or new symptomatic fibroids. Prognostic factors for better outcomes include complete removal of all visible fibroids, careful uterine reconstruction, and the absence of a strong family history of fibroids.

\section*{References}
\begin{enumerate}
  \item MedlinePlus. Myomectomy. U.S. National Library of Medicine; 2024. Available from: \url{https://medlineplus.gov/ency/article/002919.htm}
  
  \item Hoffman BL, Schorge JO, Bradshaw KD, Halvorson LM, Schaffer JI, Corton LD, editors. Williams Gynecology. 4th ed. McGraw-Hill; 2020.
  
  \item American College of Obstetricians and Gynecologists (ACOG). Management of Symptomatic Uterine Leiomyomas. Practice Bulletin No. 228. Obstet Gynecol. 2021;137(6):e100-e115.
  
  \item UpToDate. Myomectomy. Wolters Kluwer; 2024. Available from: \url{https://www.uptodate.com/contents/myomectomy}
\end{enumerate}

\clearpage